
\begin{table}[htbp]
    \centering
    \caption{Ontological Identity and Its Relation to Signature}
    \label{tab:identity-and-signature}

    \renewcommand{\arraystretch}{1.4}
    \begin{tabular}{p{3.2cm} p{5.1cm} p{5.3cm}}
        \hline
        \textbf{Conceptual Dimension} &
        \textbf{Ontological Characterization} &
        \textbf{Relation to Identity and Signature} \\
        \hline

        Identity &
        The constitutive ontological determination through which an entity is the entity that it is. Identity is inseparable from the actuality of the entity. &
        Identity is not a description, representation, or collection of identifying features. It concerns the entity itself rather than merely how it can be characterized. \\

        Signature &
        A structured configuration of features, characteristics, or determinations through which an entity may be characterized and distinguished. &
        Signature can express or indicate aspects of identity, but it must not be equated with the identity it characterizes. \\

        Identifying features &
        Particular characteristics through which an entity can be distinguished from other entities within a specified context or domain. &
        Identifying features may contribute to a signature, but not every distinguishing feature is necessarily constitutive of identity. \\

        Signature configuration &
        The organization and relation of multiple features into a recognizable or differentiating configuration. &
        A signature is more than an arbitrary list of features, but its structural organization does not establish that it exhaustively constitutes identity. \\

        Descriptive adequacy &
        The degree to which a signature or description captures relevant characteristics of an entity. &
        A descriptively adequate signature may still be ontologically incomplete. The adequacy of a description does not establish its identity with the entity's identity. \\

        Signature variability &
        The possibility that an entity may exhibit different configurations of characteristics across conditions, relations, or manifestations. &
        Variation in signature does not automatically establish a change or loss of identity. The relation between signature variation and identity preservation requires further ontological analysis. \\

        Signature continuity &
        The persistence of some signature features or configurations across a range of conditions or transformations. &
        Signature continuity may provide evidence relevant to identity preservation, but it is not sufficient by itself to establish numerical or ontological identity. \\

        Signature discontinuity &
        The alteration, disappearance, or reorganization of features previously associated with an entity. &
        Signature discontinuity does not necessarily entail identity discontinuity. Whether identity persists must be determined through an account of the entity's constitutive determinations. \\

        Signature and manifestation &
        A signature may become accessible through particular manifestations of an entity within different ontological domains. &
        A manifestation may express only some aspects of a signature. Neither manifestation nor its accessible features should be treated as exhaustive of identity. \\

        Signature and distinction &
        A signature may enable the differentiation of an entity from other entities through a configuration of distinguishing features. &
        Distinguishability through signature is not identical to the ontological constitution of the entity as itself. \\

        Signature and invariance &
        Some features of a signature may remain invariant across a specified range of transformations or variations. &
        Invariance may help characterize signature continuity, but neither signature nor identity should be reduced in advance to invariance. \\

        Non-equivalence of identity and signature &
        Identity concerns the ontological constitution of an entity, whereas signature concerns a configuration through which aspects of that entity can be characterized. &
        The relation between identity and signature is one of ontological distinction and possible expression, not conceptual equivalence or exhaustive reduction. \\

        \hline
    \end{tabular}
\end{table}
